%----------------------------------------------------------------------------------------
% Contenido del Informe - Implementación
%----------------------------------------------------------------------------------------

\chapter{Implementación}

Los códigos que se indican a continuación fueron escritos en notebooks de Colab (ver Apéndice \ref{sec:Colab}).
% \citep{Colab} permite citar a la notebook donde se escribieron los códigos.
% Ver en referencias.bib el formato con el que debe configurarse la referencia.

Para la resolución del trabajo fue necesario determinar el parámetro M, definido como el promedio de los quintos dígitos de los números de legajo de los integrantes del grupo.
\[
M=\frac{3+4+4+1}{4}=\frac{12}{4}=3
\]
En la siguiente tabla se muestra los valores de producción correspondientes al valor $M=3$.

\begin{table}[h] 
	\centering 
	\begin{tabular}{|c|c|c|}
		\toprule
		Fecha & Producción $[m^3/d]$  \\
		\midrule
		Enero 2018   & $249$ \\ \hline
		Julio 2018   & $217$ \\ \hline
		Enero 2019   & $173$ \\ \hline
		Julio 2019   & $142$ \\ \hline
		Enero 2020   & $115$ \\ \hline
		Julio 2020   & $97$ \\ \hline
		Enero 2021   & $86$ \\ \hline
		Julio 2021   & $75$  \\ \hline
		Enero 2022   & $63$  \\ \hline
		Julio 2022   & $56$  \\
		\bottomrule
	\end{tabular}
\end{table}

\section{Paquetes Utilizados}
Para empezar, se indican los paquetes que serán necesarios importar para la ejecución de los códigos.

\begin{lstlisting}[language=Python,frame=single,caption=Importación de paquetes,label={cod:paquete}]
import numpy as np
import pandas as pd
from scipy import linalg
from matplotlib import pyplot as plt

import inspect
from datetime import datetime , timedelta
\end{lstlisting}




\section{Código del cálculo de t para una determinada fecha}

Cargamos el cálculo de t para una determinada fecha con el siguiente código:

\begin{lstlisting}[language=Python,frame=single,caption=Calculo de t,label={cod:calct}] 
fecha = '01/09/2022' # Escribir la fecha como cadena (entre comillas) con el formato DD/MM/AAAA
# No editar las siguientes líneas
fechadt = datetime.strptime(fecha , '%d/%m/%Y')
diferencia = fechadt - datetime.strptime('31/12/2017' , '%d/%m/%Y')
t = diferencia.days
print(fecha , '--> t =' , t)
\end{lstlisting}



\section{Código Ajuste de Curva} 

Realizamos el ajuste de curva correspondiente con el siguiente código:

\begin{lstlisting}[language=Python,frame=single,caption=Ajuste de curva,label={cod:ajustecurva}] 
#Promedio de legajos = 3

# Ingresar en la siguiente lista los datos de abscisas (x / variable independiente)
xdata = [1,182,366,547,731,913,1097,1278,1462,1643]

# Ingresar en la siguiente lista los datos de ordenadas (y / variable dependiente)
ydata = [249,217,173,142,115,97,86,75,63,56]

# Definir valor de variable independiente cuya estimación se desea calcular
xest = 1705

# No editar las siguientes líneas
assert len(xdata)==len(ydata), 'Los vectores x e y deben ser de la misma longitud'

xlin = np.array(xdata)
ylin = np.log(ydata)
coeflin = np.polyfit(xlin,ylin,1)

A = np.exp(coeflin[1])
B = coeflin[0]

print("Ecuación de ajuste:")
ecuacion = 'f(x) = ' + str(A) + ' * e ^ (' + str(B) + ' * x)'
print(ecuacion)

xcurva = np.linspace(min(xdata),max(xdata),50)
ycurva = A * np.exp(B * xcurva)
yajuste = A * np.exp(B * np.array(xdata))
yest = A * np.exp(B * xest)

error = np.absolute(ydata - yajuste)
ecuad = error ** 2
ecm = np.mean(ecuad)
print("\nError cuadrático medio =" , ecm)

print("\nEstimación para x =" , xest)
print("f(" , xest , ") =" , yest , "\n")

plt.scatter(xdata,ydata,color='k',label='Datos')
plt.plot(xcurva,ycurva,'-r',label='Curva de ajuste')
plt.title('Ajuste exponencial')
plt.legend()
plt.show()

tabla = pd.DataFrame({"$x_i$ (datos)":xdata,"$y_i$ (datos)":ydata,"$f(x_i)$ (ajustes)":yajuste,"Error = |$y_i$ - $f(x_i)$|":error,"$Error^2$":ecuad})
print(tabla.to_latex(index=False,column_format='ccccc'))
tabla = tabla.rename(columns={"$x_i$ (datos)":"xi (datos)","$y_i$ (datos)":"yi (datos)","$f(x_i)$ (ajustes)":"f(xi) (ajustes)","Error = |$y_i$ - $f(x_i)$|":"Error = |yi - f(xi)|","$Error^2$":"Error^2"})
tabla.style.hide(axis=0)
\end{lstlisting}




\section{Códigos del Punto 1}

\subsection{Fecha de cierre del pozo}
A partir de los siguientes códigos se puede saber hasta que fecha se puede mantener el pozo abierto si consideramos que el caudal mínimo para compensar los costos de producción es de $4\frac{m^{3}}{d}$

\begin{lstlisting}[language=Python,frame=single,caption=Cálculo de cantidad de días de operación del pozo,label={cod:cantdiascierre}] 
# Q = a * e^(b * t)
# Q = 4 (m^3)/d

tcierre = np.floor(np.log(4/A) / B )
print("La cantidad de días de operacion del pozo son:", tcierre)

tcierre # No editar esta línea
\end{lstlisting}

\begin{lstlisting}[language=Python,frame=single,caption=Calculo de fecha de cierre a partir del valor de días calculado,label={cod:fechacierre}] 
fechacierre = datetime.strptime('31/12/2017' , '%d/%m/%Y') + timedelta(days=tcierre)
print('Fecha de cierre:',fechacierre.strftime('%d/%m/%Y'))
yearcierre = int(fechacierre.strftime('%Y'))
print('Año de cierre (variable yearcierre, puede usarse en gráfico del Punto 5):',yearcierre)
\end{lstlisting}



\newpage
\newpage

\section{Códigos del Punto 2} 

\subsection{Cálculo de la producción acumulada remanente, producción acumulada extraída, producción acumulada total y porcentaje remanente}

A partir del siguiente código se muestran las operaciones necesarias para completar el punto 2.

\begin{lstlisting}[language=Python,frame=single,caption=Cálculos de producción,label={cod:punto2}] 
# Escribir en la siguiente línea el cálculo para determinar la producción acumulada remanente
remanente = A * (np.exp(B*tcierre) - np.exp(B*1643)) / B

# No editar siguiente línea
print("Producción acumulada remanente =" , remanente , 'm3')

# Escribir en la siguiente línea el cálculo para determinar la producción acumulada extraída
extraida = A * (np.exp(B*1643) - np.exp(B*1)) / B

# No editar siguiente línea
print("Producción acumulada extraída =" , extraida , 'm3')

# Escribir en la siguiente línea el cálculo para determinar la producción acumulada total
total = A * (np.exp(B*tcierre) - np.exp(B*1)) / B

# No editar siguiente línea
print("Producción acumulada total =" , total , 'm3')

# Escribir en la siguiente línea el cálculo para determinar el porcentaje remanente
porc = (remanente / total) * 100

# No editar siguiente línea
print("Porcentaje de producción acumulada remanente =" , porc , '%')
\end{lstlisting}

\subsection{Código del gráfico circular} 

Para obtener el gráfico circular (de torta) con la producción acumulada remanente y la producción acumulada extraída se utilizó el siguiente código.

\begin{lstlisting}[language=Python,frame=single,caption=Gráfico de torta,label={cod:graftorta}] 
etiquetas = ['Extraída','Remanente']
volumenes = [extraida,remanente]

plt.pie(volumenes,autopct='%.1f%%',labels=etiquetas)
plt.title('Producción acumulada')
plt.show()
\end{lstlisting}

\section{Códigos del Punto 3} 

\subsection{Cálculos para la producción acumulada y días correspondientes al 50\% del volumen total}

Con el siguiente código se obtiene la producción acumulada correspondiente al 50\% del total y en cuántos días desde el inicio de la producción se llega a ese volumen.

\begin{lstlisting}[language=Python,frame=single,caption=Producción acumulada y días al 50\% del total ,label={cod:diasvol50}] 
vol50 = total/2  
print("El valor del 50% del volumen total es:", vol50)

t50 = np.ceil((np.log(vol50*(B/A) + np.exp(B*1))) / B) 

print("La cantidad de días para alcanzar el 50% del volumen total es:", t50)
t50
\end{lstlisting}

\subsection{Fecha del 50\% del volumen total}

Con el siguiente código se calcula la fecha en la que se alcanza el 50\% del volumen total a partir del valor calculado en \ref{cod:diasvol50}.

\begin{lstlisting}[language=Python,frame=single,caption=Fecha del 50\% del volumen total ,label={cod:fechavol50}] 
fechacierre = datetime.strptime('31/12/2017' , '%d/%m/%Y') + timedelta(days=t50)
print('Fecha para alcanzar el 50% del volumen total:',fechacierre.strftime('%d/%m/%Y'))
\end{lstlisting}


\section{Código del Punto 4} 

\subsection{Cálculo del valor de la producción acumulada remanente}

Para saber cual es el valor en US\$ de la producción acumulada remanente se usó el código:

\begin{lstlisting}[language=Python,frame=single,caption=Valor de la producción acumulada remanente ,label={cod:valremanente}] 
#-----------------------------------------------
# Valor de la producción acumulada remanente 
#-----------------------------------------------
# Escribir en la siguiente línea el cálculo del valor en US$ de la producción acumulada remanente
volbarriles = R * 6.28981 # Cantidad de barriles
valorrem = volbarriles * 62
print('Valor de producción acumulada remanente = US$',valorrem) # No editar esta líena
\end{lstlisting}

\newpage

\section{Código del Punto 5}

\subsection{Cálculo de la producción acumulada en un determinado año}

A partir del siguiente código se calcula la producción acumulada correspondiente a cada año, considerando a $t$ como los días de inicio y fin de cada período. Para el año 2030 se utiliza la fecha exacta de cierre del pozo calculada en el Punto 1.

\begin{lstlisting}[language=Python,frame=single,caption=Producción acumulada en una determinada temporada,label={cod:punto5}]
# --- PUNTO 5: CÁLCULO ANUAL (GRUPO 4) ---

t1ene = 1
t31dic = 365
acumyear = ((A / B) * (np.exp(B * t31dic) - np.exp(B * t1ene))) * 6.28981 * 62
print("Producción del 2018 en US$: ", np.ceil(acumyear))

t1ene = 366
t31dic = 730
acumyear = ((A / B) * (np.exp(B * t31dic) - np.exp(B * t1ene))) * 6.28981 * 62
print("Producción del 2019 en US$: ", np.ceil(acumyear))

t1ene = 731
t31dic = 1096
acumyear = ((A / B) * (np.exp(B * t31dic) - np.exp(B * t1ene))) * 6.28981 * 62
print("Producción del 2020 en US$: ", np.ceil(acumyear))

t1ene = 1097
t31dic = 1461
acumyear = ((A / B) * (np.exp(B * t31dic) - np.exp(B * t1ene))) * 6.28981 * 62
print("Producción del 2021 en US$: ", np.ceil(acumyear))

t1ene = 1462
t31dic = 1826
acumyear = ((A / B) * (np.exp(B * t31dic) - np.exp(B * t1ene))) * 6.28981 * 62
print("Producción del 2022 en US$: ", np.ceil(acumyear))

t1ene = 1827
t31dic = 2191
acumyear = ((A / B) * (np.exp(B * t31dic) - np.exp(B * t1ene))) * 6.28981 * 62
print("Producción del 2023 en US$: ", np.ceil(acumyear))

t1ene = 2192
t31dic = 2557
acumyear = ((A / B) * (np.exp(B * t31dic) - np.exp(B * t1ene))) * 6.28981 * 62
print("Producción del 2024 en US$: ", np.ceil(acumyear))

t1ene = 2558
t31dic = 2922
acumyear = ((A / B) * (np.exp(B * t31dic) - np.exp(B * t1ene))) * 6.28981 * 62
print("Producción del 2025 en US$: ", np.ceil(acumyear))

t1ene = 2923
t31dic = 3287
acumyear = ((A / B) * (np.exp(B * t31dic) - np.exp(B * t1ene))) * 6.28981 * 62
print("Producción del 2026 en US$: ", np.ceil(acumyear))

t1ene = 3288
t31dic = 3652
acumyear = ((A / B) * (np.exp(B * t31dic) - np.exp(B * t1ene))) * 6.28981 * 62
print("Producción del 2027 en US$: ", np.ceil(acumyear))

t1ene = 3653
t31dic = 4018
acumyear = ((A / B) * (np.exp(B * t31dic) - np.exp(B * t1ene))) * 6.28981 * 62
print("Producción del 2028 en US$: ", np.ceil(acumyear))

t1ene = 4019
t31dic = 4383
acumyear = ((A / B) * (np.exp(B * t31dic) - np.exp(B * t1ene))) * 6.28981 * 62
print("Producción del 2029 en US$: ", np.ceil(acumyear))

# Producción de 2030 hasta la fecha exacta de cierre
t1ene = 4384
t31dic = int(tcierre)

acumyear = ((A / B) * (np.exp(B * t31dic) - np.exp(B * t1ene))) * 6.28981 * 62
print("Producción del 2030 en US$: ", np.ceil(acumyear))

acumyear # No editar esta línea
\end{lstlisting}

\subsection{Gráfico de barras}
A partir del siguiente código se obtiene el gráfico de barras de los valores anuales de producción desde 2022 hasta el cierre del pozo en 2030.

\begin{lstlisting}[language=Python,frame=single,caption=Gráfico de barras,label={cod:grafbarras}]
# Lista de años y sus valores correspondientes
year = [2022, 2023, 2024, 2025, 2026, 2027, 2028, 2029, 2030]

# Valores de producción anual en US$
valores = [
    7514434,
    5356702,
    3827368,
    2719551,
    1938646,
    1381974,
    987421,
    701616,
    63573
]

# Gráfico de barras
plt.bar(year, valores)
plt.title('Producción Anual (US$) - 2022 al Cierre')
plt.xlabel('Año')
plt.ylabel('Valor en US$')
plt.xticks(year)
plt.show()
\end{lstlisting}
